\metatitle{Permutations}
\metadescription{Permutation matrices, Rothe diagrams, Lehmer codes, reduced words, Bruhat order and more.}




For special families of permutations, see \hyperref[namedPermutations]{the named permutations page}.


\section[permutationMatrix]{Permutation matrices}


For permutations $\pi,\sigma \in \symS_n$, we think
of multiplication as composition of functions, so that
\[
  (\pi \circ \sigma)(k) = \pi(\sigma(k)) \text{ for $k \in [n]$.}
\]
Given $\pi$, we associate its \defin{permutation matrix},
\[
 P_\pi = \left( \delta_{i, \pi(j)} \right)_{1 \leq i, j \leq n}.
\]
The permutation matrix of the composition $\pi \circ \sigma$ is $P_{\pi} P_{\sigma}$.
We follow the convention in \cite{Kitaev2011}
and illustrate permutation matrices by letting row indices start at the \emph{bottom},
as in the figure below.
From now on, this is the picture we have in mind when referring to permutation matrices.

\begin{example*}[The permutation matrix for 18537462]
The permutation matrix $P_\pi$ associated with $\pi = [1, 8, 5, 3, 7, 4, 6, 2]$
using conventional row/column coordinates, and the convention we use.
The rightmost figure is simply the graph of the function $i \mapsto \pi(i)$
(using Cartesian coordinates).
\begin{figure}
$
P_\pi =
\begin{pmatrix}
  1 & 0 & 0 & 0 & 0 & 0 & 0 & 0 \\
 0 & 0 & 0 & 0 & 0 & 0 & 0 & 1 \\
 0 & 0 & 0 & 1 & 0 & 0 & 0 & 0 \\
 0 & 0 & 0 & 0 & 0 & 1 & 0 & 0 \\
 0 & 0 & 1 & 0 & 0 & 0 & 0 & 0 \\
 0 & 0 & 0 & 0 & 0 & 0 & 1 & 0 \\
 0 & 0 & 0 & 0 & 1 & 0 & 0 & 0 \\
 0 & 1 & 0 & 0 & 0 & 0 & 0 & 0 \\
\end{pmatrix}
$
\begin{ytableau}
\none 8& \, & \bullet & \, & \, & \, & \, & \, & \, \\
\none 7& \, & \, & \, & \, & \bullet & \, & \, & \, \\
\none 6&\, & \, & \, & \, & \, & \, & \bullet & \, \\
\none 5& \, & \, & \bullet & \, & \, & \, & \, & \, \\
\none 4 & \, & \, & \, & \, & \, & \bullet & \, & \, \\
\none 3 & \, & \, & \, & \bullet & \, & \, & \, & \, \\
\none 2 & \, & \, & \, & \, & \, & \, & \, & \bullet \\
\none 1 & \bullet & \, & \, & \, & \, & \, & \, & \, \\
\none   & \none 1 &\none 2 &\none 3 &\none 4 &\none 5 &\none 6 &\none 7 &\none 8
\end{ytableau}
\end{figure}
\end{example*}

The \defin{Rothe diagram} (named after \name[H. A. Rothe]{Heinrich August Rothe}) of a permutation $\pi \in \symS_n$,
denoted $D(\pi)$, is defined as the set
\[
D(\pi) \coloneqq \{(\pi_j,i) : i < j \text{ and } \pi_i > \pi_j \}.
\]
That is, we draw the permutation matrix for $\pi$ and strike out the boxes appearing somewhere to the right or above a dot (using our convention).
This is illustrated in the figure below.
\begin{figure}
  \svgimg[width=0.25\textwidth]{svg-images/rothe-diagram.svg}{The Rothe diagram for the permutation $\pi = [1, 8, 5, 3, 7, 4, 6, 2]$.}
\end{figure}

\subsection[rotheDiagram]{Rothe diagram properties}

The number of boxes in row $i$ of $D(\pi)$ equals the number of inversions involving position $i$,
i.e., $|\{j > i : \pi_j < \pi_i\}|$.
The total number of boxes is $|D(\pi)| = \inv(\pi)$.


\subsection[lehmerCode]{Lehmer code}

The \defin{Lehmer code} (or \defin{inversion table}) of a permutation $\pi \in \symS_n$ is the sequence $c(\pi) = (c_1, c_2, \dotsc, c_n)$ where
\[
c_i \coloneqq |\{j > i : \pi_j < \pi_i\}|.
\]
Note that $0 \leq c_i \leq n - i$ for all $i$, and $c_i$ equals the number of boxes in row $i$ of $D(\pi)$.
The map $\pi \mapsto c(\pi)$ is a bijection from $\symS_n$ to the set of sequences $(c_1, \dotsc, c_n)$ with $0 \leq c_i \leq n - i$.
In particular, $\inv(\pi) = c_1 + c_2 + \dotsb + c_n$.

\begin{example*}[Lehmer code]
For $\pi = [1,8,5,3,7,4,6,2]$, we compute $c(\pi) = (0,6,3,2,3,2,1,0)$.
For instance, $c_2 = 6$ since the entries after position $2$ that are less than $\pi_2 = 8$ are $5,3,7,4,6,2$ --- all six of them.
Similarly, $c_3 = 3$ since $3,4,2$ are less than $\pi_3 = 5$ among the entries after position $3$.
Therefore $\inv(\pi) = 0+6+3+2+3+2+1+0 = 17$.
\end{example*}


\section[reducedWord]{Reduced words}

The \defin{simple transpositions} $s_1, s_2, \dotsc, s_{n-1}$
are the generators of $\symS_n$, where $s_i$ swaps positions $i$ and $i+1$.
Every permutation $\sigma \in \symS_n$ can be written as a product
$\sigma = s_{i_1} s_{i_2} \dotsm s_{i_\ell}$.
The minimum such $\ell$ is the \defin{length} of $\sigma$,
denoted $\ell(\sigma)$, and equals the number of inversions $|\INV(\sigma)|$.
A product $s_{i_1} s_{i_2} \dotsm s_{i_\ell}$ achieving this minimum
is called a \defin{reduced word} (or \defin{reduced expression}) for $\sigma$.


\subsection[reducedWordExamples]{Examples}

Every $\sigma \in \symS_n$ satisfies $\ell(\sigma) = \inv(\sigma) = |D(\sigma)|$.

\begin{example*}[Reduced words]
Consider $\sigma = [3,1,4,2]$. We have $\inv(\sigma) = 3$
(the inversions are $(1,2), (1,4), (3,4)$), so $\ell(\sigma) = 3$.
One reduced word is $s_2 s_1 s_3$:
\[
[1,2,3,4] \xrightarrow{s_2} [1,3,2,4] \xrightarrow{s_1} [3,1,2,4] \xrightarrow{s_3} [3,1,4,2].
\]
Since $s_1$ and $s_3$ commute (they act on disjoint positions),
$s_2 s_3 s_1$ is also a reduced word:
\[
[1,2,3,4] \xrightarrow{s_2} [1,3,2,4] \xrightarrow{s_3} [1,3,4,2] \xrightarrow{s_1} [3,1,4,2].
\]
These are the only two reduced words for $[3,1,4,2]$.
The number of reduced words for $\sigma$ is given by the
coefficient of $x_1 x_2 \dotsm x_n$ in the \hyperref[stanleySym]{Stanley symmetric function} $F_\sigma(\xvec)$.
\end{example*}


\subsection[bruhatOrder]{Bruhat order and weak order}\label{weakOrder}

The \defin{Bruhat order} (or \defin{strong order}) on
permutations is a partial order $\leq$ on $\symS_n$.
We have that $\sigma \leq \tau$ if
$\tau$ can obtained from $\sigma$ via a sequence of transpositions,
which increase the number of inversions in each step.
Alternatively, we have that $\sigma \leq \tau$ if and only if
there are reduced words $w_\sigma$, $w_\tau$ for $\sigma$ and $\tau$
such that $w_\sigma$ is a subword of $w_\tau$.

The Bruhat order is a ranked poset, with rank given by the number of inversions.

The (right) \defin{weak order} is a subposet of the Bruhat order,
where only \emph{simple transpositions} are allowed, when multiplying on the right.
This corresponds to swapping adjacent entries.
Let us write $\leq_R$ for this relation.
Alternatively, we have that $\sigma \leq_R \tau$ if and only if
there are reduced words $w_\sigma$, $w_\tau$ for $\sigma$ and $\tau$
such that $w_\sigma$ is a \emph{prefix} of $w_\tau$.


\begin{figure}
  \svgimg[width=0.5\textwidth]{svg-images/bruhatGraph.svg}{Bruhat poset for $n=4$.}
\end{figure}

In this figure, we show the Hasse diagram of the Bruhat order for $\symS_4$.
The solid edges indicate covering relations in the (right) weak order.
See \cite[Ch. 10.5]{Fulton1997}, \cite[Ex. 183, 185]{StanleyEC1}
and \cite[Ch.2-Ch.3]{Bjorner2005} for more properties
and alternative definitions.


\subsection[braidRelations]{Braid relations}


The symmetric group $\symS_n$ is generated by the simple transpositions
$s_1, s_2, \dotsc, s_{n-1}$, and these generators satisfy the relations
\[
s_i^2 = 1,\qquad s_i s_j = s_j s_i \text{ if } |i-j| > 1,\qquad
s_i s_{i+1} s_i = s_{i+1} s_i s_{i+1}.
\]
The last relation is called the \defin{braid relation}, and together these
relations give the usual Coxeter presentation of $\symS_n$.
They also explain how different reduced words for the same permutation
arise: any two reduced words are connected by a sequence of commutation
relations and braid relations.




For \hyperref[notationAffinePermutations]{affine permutations},
\hyperref[permutationsTypeB]{signed permutations},
\hyperref[permutationsColored]{colored permutations},
and \hyperref[permutationsStirling]{Stirling permutations},
see the \hyperref[notationAffinePermutations]{generalizations of permutations} page.
